\section{Proposal and Positioning}
\label{sec:proposal}

\subsection{The Problem: Redundancy in Symbolic Regression Search Spaces}

Symbolic regression (SR) seeks a mathematical expression $f(x_1, \ldots, x_m)$
that best fits observed data $(X, y)$. Modern approaches represent expressions
as Directed Acyclic Graphs (DAGs) where leaf nodes are variables or constants,
and internal nodes are operations ($+, \times, \sin, \cos$, etc.)
\cite{liu2025, petersen2021}.

A fundamental issue plagues all existing SR methods: the search space contains
\textbf{massive redundancy}. For an expression DAG with $k$ internal nodes,
there exist $O(k!)$ distinct node labelings that all represent the
\emph{same expression}. Each of these $k!$ variants occupies a separate point
in the search space, wasting computational effort on equivalent candidates.

\begin{proposition}[Search space redundancy]
For an expression DAG $D$ with $k$ internal (non-variable) nodes,
the number of structurally equivalent but differently-labeled DAGs is
$\Theta(k!)$ in the worst case.
\end{proposition}

\subsection{Our Contribution: Isomorphism-Invariant Representation}

IsalSR extends the IsalGraph framework \cite{lopezrubio2025}---originally
designed for topology-only unlabeled graphs---to \textbf{labeled DAGs} for
symbolic regression. The key contribution is the \textbf{canonical string}
$w^*_D$: a complete labeled-DAG invariant that satisfies
\begin{equation}
  w^*_{D_1} = w^*_{D_2} \iff D_1 \cong D_2
  \label{eq:invariant}
\end{equation}
where $\cong$ denotes labeled-DAG isomorphism (preserving node labels and
edge directions).

By mapping every expression to its canonical string before evaluation,
the search space is reduced from $O(N \cdot k!)$ to $O(N)$ unique expressions,
where $N$ is the number of structurally distinct DAGs. This is the
\textbf{sole contribution}: we provide an invariant representation that
existing SR methods can plug into for acceleration.

\subsection{What IsalSR Is NOT}

Following the advisor's explicit guidance, we emphasize:
\begin{itemize}[nosep]
  \item IsalSR is \textbf{not a new SR method}. It is a representation layer.
  \item We do \textbf{not claim SR expertise}. Our expertise is in graph theory
    and combinatorics applied to expression representation.
  \item The search operators (mutation, crossover, hill climbing, evolutionary
    population) are \textbf{standard, generic wrappers}---not novel algorithms.
  \item The experimental design compares the \textbf{same algorithm} operating
    in canonical vs.\ non-canonical space, not different algorithms.
\end{itemize}

\subsection{Relationship to IsalGraph}

IsalGraph \cite{lopezrubio2025} represents arbitrary \emph{unlabeled, undirected}
graphs as instruction strings. IsalSR adapts this framework with four key
modifications:
\begin{enumerate}[nosep]
  \item \textbf{Labeled nodes}: Each node has a type (VAR, ADD, MUL, SIN, etc.)
    encoded in two-character compound tokens.
  \item \textbf{Directed edges}: Edges represent data flow
    ($u \to v$ means ``$u$ provides input to $v$'').
  \item \textbf{DAG constraint}: The C/c edge instructions enforce acyclicity
    via cycle detection (silently skipping cycle-creating edges).
  \item \textbf{Fixed starting node}: The canonical string is computed from
    $x_1$ only (the first input variable), since variables are distinguishable
    and pre-numbered. This eliminates the $O(n)$ loop over starting nodes
    present in IsalGraph.
\end{enumerate}
